\chapter{\abstractname}

%TODO: Abstract

Cardiovascular diseases and cancer cause 5.2 million deaths in Europe and 27.4 million deaths worldwide every year. Minimally invasive endovascular interventions play important role in the treatment of the above diseases, because they reduce the risk of surgery and shorten the hospitalization of patients. These interventions require state-of-the-art X-ray angiography technology. (“X-ray angiography” refers to real-time X-ray imaging of the interior of arteries or veins following the injection of contrast media. The procedure allows the interventional radiologist to diagnose and treat problems within the blood vessel.) The currently used Digital Subtraction Angiograph (\gls{dsa}) technology has several limitations and drawbacks.

Kinepict Health Ltd. is a Hungarian startup formed as a spin-off of the Semmelweis University, Budapest. In response to the challenges faced by DSA, the Kinepict team developed and clinically validated Digital Variance Angiography (\gls{dva}). DVA uses advanced statistical and mathematical tools to extract movement related functional information from the X-ray image series. As other imaging techniques, DVA requires further postprocessing in order to improve visibility of the vessel region. This thesis introduces, enhances and combines such algorithms.



\makeatletter
\ifthenelse{\pdf@strcmp{\languagename}{english}=0}
{\renewcommand{\abstractname}{Kurzfassung}}
{\renewcommand{\abstractname}{Abstract}}
\makeatother

%\chapter{\abstractname}

%TODO: Abstract in other language
\begin{otherlanguage}{ngerman} % TODO: select other language, either ngerman or english !

\end{otherlanguage}


% Undo the name switch
\makeatletter
\ifthenelse{\pdf@strcmp{\languagename}{english}=0}
{\renewcommand{\abstractname}{Abstract}}
{\renewcommand{\abstractname}{Kurzfassung}}
\makeatother